\section{All cubic Tornheim directions from two scalar coordinates}
\label{tr3:section}

The incoming \emph{Coincident Stieltjes and Directional Zeta} report~\cite{tr3:incoming} proves
all quadratic derivatives of $T(at,bt,ct)$ and asks for the entire cubic
layer.  The following calculation answers that question with two explicit
scalar coordinates.  One is the already isolated diagonal derivative
$\Omega$; the other is an absolutely convergent logarithmic Gamma series.
The result also identifies an exact conic of directions and a cyclic
symmetrization in which the second coordinate cancels.  These are finite
identities; no arithmetic independence of the two coordinates is assumed.

Throughout this section,
\[
 T(A,B,C)=\sum_{m,n\geq1}m^{-A}n^{-B}(m+n)^{-C},
 \qquad L=\log(2\pi),
\]
initially in its domain of absolute convergence and subsequently by its
meromorphic continuation.  Set
\[
 z_2=\zeta''(0),\qquad z_3=\zeta'''(0),\qquad
 y_3=\zeta'''(-1),\qquad
 \Omega=\left.\frac{d^3}{dt^3}T(t,t,t)\right|_{t=0}.
\]
The order of restriction is part of the notation.  In particular, $T$ is
not jointly holomorphic at the origin, whereas the one-variable diagonal
is holomorphic there~\cite{tr3:BD}.  The Stieltjes convention is
\[
 \zeta(1+u)=\frac1u+\sum_{r\geq0}\frac{(-1)^r\gamma_r}{r!}u^r.
\]

\subsection{A local remainder with no Gamma factors in its polar part}

\begin{lemma}[Elementary polar model]\label{tr3:polar}
There is a unique function $R(A,B,C)$ holomorphic in a neighborhood of
$(0,0,0)$, symmetric in $A,B$, for which
\begin{equation}\label{tr3:model}
 \boxed{\displaystyle
 T(A,B,C)=\zeta(A)\zeta(B)
  -\frac{C\zeta(B-1)}{A+C}
  -\frac{C\zeta(A-1)}{B+C}+C R(A,B,C).}
\end{equation}
The identity is meromorphic; each quotient retains its indicated
numerator until a restriction has been specified.
\end{lemma}

\begin{proof}
For $A,B$ near zero, Jonqui\`ere's expansion~\cite[\S25.12(ii)]{tr3:DLMF} gives
\[
 \operatorname{Li}_A(e^{-x})
 =\Gamma(1-A)x^{A-1}
   +\sum_{j\geq0}\zeta(A-j)\frac{(-x)^j}{j!},\qquad 0<x<2\pi.
\]
In the Mellin integral
\[
 \Gamma(C)T(A,B,C)=\int_0^\infty
 x^{C-1}\operatorname{Li}_A(e^{-x})\operatorname{Li}_B(e^{-x})\,dx,
\]
subtract on $(0,1)$ the six terms
\begin{align*}
 S(A,B;x)={}&\Gamma(1-A)\Gamma(1-B)x^{A+B-2}
 +\Gamma(1-A)\zeta(B)x^{A-1}
 +\Gamma(1-B)\zeta(A)x^{B-1}\\
 &-\Gamma(1-A)\zeta(B-1)x^A
 -\Gamma(1-B)\zeta(A-1)x^B+\zeta(A)\zeta(B).
\end{align*}
The difference is normally integrable against $x^{C-1}$ on a sufficiently
small common parameter polydisc.  Each fixed parameter derivative only
adds logarithmic powers.  Integrating the six displayed monomials shows
that some holomorphic symmetric $H$ satisfies
\begin{align}\label{tr3:Hmodel}
 T(A,B,C)=\frac{1}{\Gamma(1+C)}\biggl\{
 \zeta(A)\zeta(B)
 -\frac{C\Gamma(1-A)\zeta(B-1)}{A+C}
 -\frac{C\Gamma(1-B)\zeta(A-1)}{B+C}+C H(A,B,C)\biggr\}.
\end{align}
The terms with denominators $A+B+C-2$, $A+C-1$, and $B+C-1$ belong
to $H$, since these denominators are nonzero at the origin.

Write $G(C)=1/\Gamma(1+C)$.  Subtract the proposed polar part from
\eqref{tr3:Hmodel} and divide by $C$.  The resulting expression is
\begin{align*}
 R(A,B,C)={}&G(C)H(A,B,C)
  +\frac{G(C)-1}{C}\zeta(A)\zeta(B)\\
 &-\frac{\Gamma(1-A)G(C)-1}{A+C}\zeta(B-1)
  -\frac{\Gamma(1-B)G(C)-1}{B+C}\zeta(A-1).
\end{align*}
Every quotient is holomorphic.  For example,
$\Gamma(1-A)G(-A)=1$, so its numerator is divisible by $A+C$ in
the local ring of holomorphic functions.  This proves existence and
symmetry.  Uniqueness follows from \eqref{tr3:model} for $C\ne0$ and
then by analytic continuation.
\end{proof}

The meromorphic continuation and Mellin method are classical~\cite{tr3:BD,tr3:SS}; the local
form \eqref{tr3:Hmodel} also occurs in the incoming report.  The change
from $H$ to $R$ is useful because all Gamma and Euler constant terms
are absorbed into a holomorphic function before the coefficient
calculation.

For reference, the complete linear jet is
\begin{align}\label{tr3:linearjet}
 R(0,0,0)&=\frac L2+\zeta'(-1),\notag\\
 R_A(0)=R_B(0)&=\frac{L^2}{8}-\frac{z_2}{2},\qquad
 R_C(0)=\frac{\zeta''(-1)-z_2}{2}.
\end{align}
These formulas need no additional spectral input.  The exact $C$ axis
in the next proof determines $R(0)$ and $R_C(0)$.  The coefficient of
$AC$ in the Euler identity \eqref{tr3:stuffleR} gives
$2R_A=\zeta'(0)^2-\zeta''(0)$.  As a consequence,
\[
 \left.\frac{d^2}{dt^2}T(t,t,t)\right|_{t=0}
  =\frac{L^2}{2}-z_2-\zeta''(-1)+4R_A+2R_C
  =L^2-4z_2.
\]
This recovers the incoming quadratic diagonal formula by an independent
coefficient calculation, while the new result begins at cubic order.

\subsection{The quadratic remainder and the cubic ray theorem}

Let all partial derivatives in the following notation be evaluated
at $(0,0,0)$:
\[
 p=R_{AA}=R_{BB},\qquad q=R_{AB},\qquad
 r=R_{AC}=R_{BC},\qquad s=R_{CC}.
\]
These four coordinates describe the homogeneous quadratic part of
$R$.  Three exact restrictions determine three independent combinations.

\begin{lemma}[Three exact constraints]\label{tr3:constraints}
The four second partial derivatives satisfy
\begin{equation}\label{tr3:quadraticjet}
 \boxed{\displaystyle
 s=\frac{y_3-z_3}{3},\qquad
 p+2r=-\frac{Lz_2}{2}-z_3,\qquad
 q=\frac{\Omega+6Lz_2+8z_3}{6}.}
\end{equation}
\end{lemma}

\begin{proof}
Counting the $n-1$ positive pairs with sum $n$ gives
\[
 T(0,0,C)=\zeta(C-1)-\zeta(C).
\]
Since $\zeta(0)^2-2\zeta(-1)=5/12$, \eqref{tr3:model} yields
\[
 C R(0,0,C)=\zeta(C-1)-\zeta(C)-\frac5{12}.
\]
Its cubic coefficient gives $s=(y_3-z_3)/3$.

On the plane $B=0$, set
\[
 E(A,C)=T(A,0,C)=\sum_{n>m\geq1}n^{-C}m^{-A}.
\]
The ordinary Euler stuffle identity, first in an absolutely convergent
region and then meromorphically, is
\[
 E(A,C)+E(C,A)=\zeta(A)\zeta(C)-\zeta(A+C).
\]
Inserting \eqref{tr3:model} gives a holomorphic identity
\begin{align}\label{tr3:stuffleR}
 C R(A,0,C)+A R(C,0,A)
 ={}&\zeta(A)\zeta(C)-\zeta(A+C)
 +\frac{\zeta(A)+\zeta(C)}2\notag\\
 &+\zeta(-1)+\zeta(A-1)+\zeta(C-1).
\end{align}
The coefficient of $A^2C$ on its left is $p/2+r$.  On its right it is
$\zeta'(0)z_2/2-z_3/2=-Lz_2/4-z_3/2$.  This proves the second
constraint.

Finally, on the diagonal \eqref{tr3:model} says
\[
 T(t,t,t)=\zeta(t)^2-\zeta(t-1)+tR(t,t,t).
\]
Taking third derivatives gives
\[
 \Omega=-3Lz_2-z_3-y_3+6p+6q+12r+3s.
\]
The first two constraints simplify this to
$\Omega=-6Lz_2-8z_3+6q$, proving the last assertion.
\end{proof}

\begin{theorem}[Every asymmetric cubic direction]\label{tr3:main}
Let $a,b,c\in\mathbb C$ satisfy $(a+c)(b+c)\ne0$.  Form the
one-variable meromorphic restriction $W_{a,b,c}(t)=T(at,bt,ct)$ before
continuing to $t=0$.  This restriction is holomorphic at zero.  Define
$\kappa=R_{AA}(0,0,0)$.  Then
\begin{align}\label{tr3:allrays}
 \boxed{\displaystyle W_{a,b,c}'''(0)}
 ={}&abc\,\Omega
 +3c\bigl(a^2+b^2-c(a+b)\bigr)\kappa\notag\\
 &-\frac32\bigl(ab(a+b-4c)+c^2(a+b)\bigr)Lz_2\notag\\
 &+\left(-\frac{a^3+b^3}{2}+8abc-3c^2(a+b)-c^3\right)z_3\notag\\
 &+\left(c^3-\frac{cb^3}{a+c}-\frac{ca^3}{b+c}\right)y_3.
\end{align}
The scalar $\kappa$ is explicitly given by the absolutely convergent
series in \eqref{tr3:kappa} below.
\end{theorem}

\begin{proof}
On an admissible ray the two polar quotients in \eqref{tr3:model}
become constants times holomorphic one-variable zeta functions, so
the restriction is holomorphic.  Its cubic derivative is
\begin{align*}
 W_{a,b,c}'''(0)={}&-\frac{a^3+b^3}{2}z_3
 -\frac32ab(a+b)Lz_2
 -c\left(\frac{b^3}{a+c}+\frac{a^3}{b+c}\right)y_3\\
 &+3cp(a^2+b^2)+6cqab+6c^2r(a+b)+3c^3s.
\end{align*}
Substitute $p=\kappa$ and \eqref{tr3:quadraticjet}.  Collecting the
coefficients of $\Omega,\kappa,Lz_2,z_3,y_3$ gives
\eqref{tr3:allrays}.
\end{proof}

Thus no unresolved partial derivative of the three-variable Tornheim
function remains in the answer.  The new information relative to the
quadratic theorem is one explicitly specified Gamma-series coordinate
in addition to the diagonal coordinate already present in the cubic
digamma problem.  This is a two-coordinate reduction, not a claim that
two algebraically independent constants are necessary.

\subsection{A convergent Gamma series and all boundary jets}

Put
\begin{equation}\label{tr3:rho}
 \rho(n)=\log\Gamma(n+1)-\left(n+\frac12\right)\log n
            +n-\frac L2-\frac{1}{12n},\qquad n\geq1.
\end{equation}
Stirling's expansion~\cite[\S5.11]{tr3:DLMF} gives $\rho(n)=-1/(360n^3)+O(n^{-5})$.
Consequently
\[
 M_j=\sum_{n\geq1}(\log n)^j\rho(n)
\]
converges absolutely for every fixed integer $j\geq0$.  No zeta
regularization is used in these sums.

\begin{theorem}[Boundary germ in Gamma and Stieltjes data]
\label{tr3:boundary}
The Dirichlet series
\[
 D(A)=\sum_{n\geq1}\frac{\rho(n)}{n^A}
\]
is holomorphic on $\Re A>-2$.  Near $A=0$ it satisfies
\begin{align}\label{tr3:boundarygerm}
 R(A,0,0)={}&D(A)-\zeta'(A-1)-\frac12\zeta'(A)
            -\zeta(A-1)
            +\frac{\zeta(1+A)-1/A}{12}.
\end{align}
In particular, for every $j\geq0$,
\begin{align}\label{tr3:boundaryjets}
 \partial_A^jR(0,0,0)
 ={}&(-1)^j M_j-\zeta^{(j+1)}(-1)
 -\frac12\zeta^{(j+1)}(0)-\zeta^{(j)}(-1)
 +\frac{(-1)^j\gamma_j}{12},
\end{align}
and the scalar in Theorem~\ref{tr3:main} is
\begin{equation}\label{tr3:kappa}
 \boxed{\displaystyle
 \kappa=\sum_{n\geq1}(\log n)^2\rho(n)
       -\zeta'''(-1)-\frac12\zeta'''(0)-\zeta''(-1)
       +\frac{\gamma_2}{12}.}
\end{equation}
\end{theorem}

\begin{proof}
The decay of $\rho$ proves normal convergence of $D$ and every fixed
order derivative on compact subsets of $\Re A>-2$.  For $\Re A>2$
and $C$ near zero, the representation
\[
 E(A,C)=\sum_{m\geq1}m^{-A}\zeta(C,m+1)
\]
can be differentiated normally in $C$.  Lerch's formula
$\zeta'(0,x)=\log\Gamma(x)-L/2$~\cite[(25.11.18)]{tr3:DLMF} gives
\[
 E_C(A,0)=\sum_{m\geq1}m^{-A}
            \left(\log\Gamma(m+1)-\frac L2\right).
\]
Expanding the summand using \eqref{tr3:rho}, in this convergent
half-plane, proves
\[
 E_C(A,0)=D(A)-\zeta'(A-1)-\frac12\zeta'(A)
          -\zeta(A-1)+\frac{\zeta(A+1)}{12}.
\]
Both sides have meromorphic continuations to a neighborhood of zero.
On the other hand, differentiation in $C$ of \eqref{tr3:model} with
$B=0$, for fixed $A\ne0$, gives
\[
 E_C(A,0)=\frac{1}{12A}+R(A,0,0).
\]
Combining these two identities proves \eqref{tr3:boundarygerm}; its
apparent $A=0$ singularity is removable.  Taking $j$ derivatives and
using the specified Stieltjes convention proves
\eqref{tr3:boundaryjets} and \eqref{tr3:kappa}.
\end{proof}

The first two members of this hierarchy are completely evaluable:
\begin{align}\label{tr3:lowGamma}
 \sum_{n\geq1}\rho(n)
   &=2\zeta'(-1)+\frac L4-\frac{1+\gamma}{12},\notag\\
 \sum_{n\geq1}(\log n)\rho(n)
   &=-\frac{L^2}{8}-\zeta''(-1)-\zeta'(-1)-\frac{\gamma_1}{12}.
\end{align}
They follow from \eqref{tr3:linearjet} and
\eqref{tr3:boundaryjets}.  We use these as derived consistency
identities and make no separate novelty claim for them.

There is also an exact integral involving an order derivative of the
ordinary polylogarithm.  In the half-plane $\Re A>-2$,
\begin{equation}\label{tr3:Dpolylog}
 D(A)=\int_0^\infty
 \left(\frac1{e^x-1}-\frac1x+\frac12-\frac{x}{12}\right)
           \frac{\operatorname{Li}_A(e^{-x})}{x}\,dx.
\end{equation}
To see this, the Laplace form of Binet's formula~\cite[\S5.9]{tr3:DLMF} gives
\[
 \rho(n)=\int_0^\infty e^{-nx}
 \left(\frac1{e^x-1}-\frac1x+\frac12-\frac{x}{12}\right)\frac{dx}{x}.
\]
This formula can also be verified by differentiating with respect to
$n$ using the standard digamma integral and fixing the constant at
$n=+\infty$.  Summation and integration commute absolutely on
$\Re A>-2$.  At the lower endpoint the bracket divided by $x$ is
$-x^2/720+O(x^4)$ and the polylogarithm contributes at worst the
integrable singularity required by that half-plane; the upper endpoint
is exponentially decaying.  Fixed order derivatives add logarithmic
factors and remain integrable.  Hence
\begin{equation}\label{tr3:Mpolylog}
 M_j=(-1)^j\int_0^\infty
 \left(\frac1{e^x-1}-\frac1x+\frac12-\frac{x}{12}\right)
 \left.\partial_A^j\operatorname{Li}_A(e^{-x})\right|_{A=0}\frac{dx}{x}.
\end{equation}
Equations \eqref{tr3:kappa} and \eqref{tr3:Mpolylog} supply equivalent
Gamma-series and polylogarithmic-integral coordinates for the same
constant.  They also connect the entire boundary derivative hierarchy
to the Stieltjes constants.  Since
$-\left.\partial_s H_n^{(s)}\right|_{s=0}=\log\Gamma(n+1)$,
these identities can equivalently be written with order derivatives of
generalized harmonic numbers.

\subsection{Exact cancellation families and examples}

\begin{corollary}[A conic of directions]\label{tr3:conic}
If an admissible direction obeys
\begin{equation}\label{tr3:coniceq}
 a^2+b^2=c(a+b),
\end{equation}
then its third derivative is a finite combination of
$\Omega$, $L\zeta''(0)$, $\zeta'''(0)$, and $\zeta'''(-1)$.
No Gamma-series coordinate occurs.
\end{corollary}

For example, the direction $(1,2,5/3)$ gives the exact relation
\begin{equation}\label{tr3:conicexample}
 W_{1,2,5/3}'''(0)=\frac{10}{3}\Omega
 -\frac32Lz_2-\frac{403}{54}z_3-\frac{245}{297}y_3.
\end{equation}
The diagonal and the Euler diagonal are familiar points on the same
conic.  In particular,
\[
 W_{1,0,1}'''(0)=-\frac32Lz_2-\frac92z_3,
\]
as also follows directly from
$T(t,0,t)=(\zeta(t)^2-\zeta(2t))/2$.

\begin{corollary}[Cyclic elimination]\label{tr3:cyclic}
Let all pairwise sums $a+b$, $a+c$, $b+c$ be nonzero, and put
$e_1=a+b+c$, $e_2=ab+bc+ca$, $e_3=abc$.  Then
\begin{align}\label{tr3:cyclicformula}
 W_{a,b,c}'''(0)+W_{b,c,a}'''(0)+W_{c,a,b}'''(0)
 ={}&3e_3\Omega+3(9e_3-e_1e_2)Lz_2\notag\\
 &+(-2e_1^3+3e_1e_2+27e_3)z_3.
\end{align}
Both $\kappa$ and $\zeta'''(-1)$ cancel.
\end{corollary}

\begin{proof}
Sum \eqref{tr3:allrays} cyclically.  The $\kappa$ coefficient is
$3\sum_{\rm cyc}c(a^2+b^2)-3\sum_{\rm cyc}c^2(a+b)=0$.
In the $y_3$ coefficient, terms sharing the denominator $a+c$ combine
to $-b^3$, cancelling the sum of cubes.  The remaining two
coefficients reduce to the elementary symmetric expressions displayed
in \eqref{tr3:cyclicformula}.
\end{proof}

The concrete specialization $(a,b,c)=(1,2,3)$ is
\[
 W_{1,2,3}'''(0)+W_{2,3,1}'''(0)+W_{3,1,2}'''(0)-18\Omega
       =-36Lz_2-72z_3.
\]
Individually, two asymmetric cases read
\begin{align*}
 W_{1,2,3}'''(0)&=6\Omega-36\kappa
                -\frac{27}{2}Lz_2-\frac{129}{2}z_3+\frac{102}{5}y_3,\\
 W_{2,3,1}'''(0)&=6\Omega+24\kappa
                -\frac{33}{2}Lz_2+\frac{29}{2}z_3-10y_3.
\end{align*}
The Euler direction $(1,0,2)$ isolates the additional coordinate:
\begin{equation}\label{tr3:kappaisolate}
 W_{1,0,2}'''(0)=-6\kappa-6Lz_2-\frac{41}{2}z_3+7y_3.
\end{equation}
Thus one may replace $\kappa$ by a single Euler double-zeta directional
derivative if that basis is preferable.  The Gamma series gives a
convergent evaluation of that derivative, rather than merely naming it.

For a finite linear combination of admissible rays, the map
\[
 (a,b,c)\longmapsto
 \bigl(abc,\ 3c(a^2+b^2-c(a+b))\bigr)
\]
lists the coefficients of the two residual coordinates.  Any linear
combination in the kernel of this explicit map reduces to the displayed
ordinary zeta derivatives.  This is a sufficient exact elimination
criterion.  Arithmetic relations between $\Omega$ and $\kappa$, if
found later, could enlarge the actual reducible family.

\subsection{Verification and status}

The accompanying \texttt{check\_cubic\_rays.py} checks the finite algebra
symbolically.  It computes the Gamma series using finite Stirling
subtraction and evaluates Tornheim functions independently through
integrated Jonqui\`ere series on $(0,1)$ and an incomplete-Gamma series
on $(1,\infty)$.  The reported working precision, finite-difference
step, truncations, and discrepancies are recorded in
\texttt{cubic\_rays.json}.  The largest discrepancy in the three fully asymmetric cases is less than
$2.27\times10^{-38}$ at 65-digit working precision.  These numerical
diagnostics accompany the proof; the floating-point results are not interval certificates.
For orientation,
\begin{align*}
 \kappa&=3.6336347857639257352355574698800935822\ldots,\\
 M_2&=-0.00064921643050147820901953512402746385\ldots.
\end{align*}

The exact contribution is the cubic directional classification,
the explicit boundary germ, and their cancellation consequences.
The classical ingredients are Mellin continuation, Euler stuffle,
Stirling subtraction, and Lerch's formula.  The 2025 preprint of
Sathyanarayana and Sharan~\cite{tr3:SS}, published in 2026, develops fixed-variable
Kronecker limits, Herglotz-type relations, and diagonal special values.
Its Section~5 proposes further study of simultaneously varying order
limits at points of indeterminacy.  The 2026 preprint of
Dobrowolski~\cite{tr3:Dobrowolski} concerns positive-integer symmetrized weights.  These are
relevant adjacent developments.  The simultaneous cubic-origin formula
above was not located in the inspected versions, but a comprehensive
priority claim would require a broader literature review.

Questions about further reduction of $\Omega$ and $\kappa$, the quartic
jet, and systematic cancellation families are stated in
Section~\ref{research:section}.
